% ============================================================================
% COURS 03 — LOIS ENTRÉE-SORTIE — PARTIE B
% Source : 3-LoiES20222.docx
% ============================================================================

\section{Modèles géométriques et contraintes de position}
\label{sec:les-modeles}

Le modèle géométrique direct d'une chaîne ouverte s'obtient généralement
par une succession de relations de Chasles. L'objectif est d'exprimer le
vecteur position de l'effecteur dans la base du bâti.

Pour un point $P$ :
\[
\overrightarrow{AP}=x_P\vec x_0+y_P\vec y_0+z_P\vec z_0.
\]
Les coordonnées sont obtenues en projetant sur les vecteurs de la base :
\[
x_P=\overrightarrow{AP}\cdot\vec x_0,
\quad y_P=\overrightarrow{AP}\cdot\vec y_0,
\quad z_P=\overrightarrow{AP}\cdot\vec z_0.
\]

\subsection{Contraintes géométriques}

Une contrainte de position limite l'ensemble des positions possibles de
l'effecteur. Elle peut s'écrire sous forme vectorielle ou scalaire.

\begin{center}
\small
\begin{tabularx}{\linewidth}{>{\bfseries}p{3.3cm}X}
\toprule
Contrainte & Écriture possible \\
\midrule
Distance à un point $A$ & $\|\overrightarrow{AP}\|\leq d$ \\
Appartenance à un plan de normale $\vec n$ & $\overrightarrow{AP}\cdot\vec n=h$ \\
Appartenance à une droite & deux coordonnées ou deux projections sont imposées \\
Trajectoire paramétrée & $\mathbf X=\mathbf X(u)$ avec $u$ paramètre de trajectoire \\
Orientation imposée & produits scalaires entre les vecteurs des repères \\
\bottomrule
\end{tabularx}
\end{center}

\begin{TLThreeApp}{Position d'un robot série}
Écrire le vecteur position de l'effecteur par une relation de Chasles,
le projeter dans la base du bâti et déterminer les coordonnées
opérationnelles. Identifier ensuite les configurations articulaires
permettant d'atteindre un point imposé.
\end{TLThreeApp}

\section{Déterminer une loi en position par fermeture géométrique}
\label{sec:les-fermeture}

Dans une chaîne fermée, le graphe des liaisons comporte au moins une
boucle. L'équation vectorielle de fermeture fournit deux équations
scalaires dans un problème plan et trois dans un problème spatial.

\begin{TLThreeMethod}{Démarche générale}
\begin{enumerate}
  \item Identifier le paramètre d'entrée $q_e$ et le paramètre de sortie $q_s$.
  \item Choisir une boucle et écrire la fermeture géométrique.
  \item Exprimer chaque vecteur à l'aide des paramètres géométriques et cinématiques.
  \item Projeter dans une base adaptée.
  \item Éliminer les paramètres intermédiaires.
  \item Mettre en évidence une relation ne contenant que $q_e$, $q_s$ et les constantes géométriques.
  \item Vérifier l'homogénéité, les positions particulières et le domaine de validité.
\end{enumerate}
\end{TLThreeMethod}

\subsection{Exemple du bielle-manivelle}

On note $e=OA$ la longueur de la manivelle, $L=AB$ la longueur de la
bielle, $x=OB$ la position du piston, $\alpha$ l'angle de la manivelle et
$\beta$ l'angle de la bielle.

\TLBielleManivelleGeometry

La fermeture s'écrit :
\[
\overrightarrow{OA}+\overrightarrow{AB}+\overrightarrow{BO}=\vec0,
\]
soit :
\[
e\vec x_1+L\vec x_2-x\vec x_0=\vec0.
\]
Les projections sur $\vec x_0$ et $\vec y_0$ donnent :
\[
\left\{\begin{aligned}
e\cos\alpha+L\cos\beta-x&=0,\\
e\sin\alpha+L\sin\beta&=0.
\end{aligned}\right.
\]

\subsection{Méthode 1 : élimination par identité trigonométrique}

On isole les termes en $\beta$ :
\[
L\cos\beta=x-e\cos\alpha,
\qquad
L\sin\beta=-e\sin\alpha.
\]
En élevant au carré et en additionnant :
\[
\boxed{L^2=x^2+e^2-2ex\cos\alpha.}
\]
Cette relation constitue une loi entrée-sortie implicite entre $x$ et
$\alpha$.

Pour éliminer une distance commune, on peut diviser deux équations afin
de faire apparaître une tangente :
\[
\tan\beta=\frac{-e\sin\alpha}{x-e\cos\alpha}.
\]

\subsection{Méthode 2 : norme au carré}

Après isolement d'un vecteur de norme connue :
\[
L\vec x_2=x\vec x_0-e\vec x_1.
\]
On prend la norme au carré :
\[
L^2=\|x\vec x_0-e\vec x_1\|^2
=x^2+e^2-2ex\,\vec x_0\cdot\vec x_1,
\]
d'où le même résultat puisque
$\vec x_0\cdot\vec x_1=\cos\alpha$.

\subsection{Méthode 3 : théorème d'Al-Kashi}

Lorsque la boucle forme directement un triangle, la loi des cosinus
permet d'écrire sans projection :
\[
\boxed{L^2=x^2+e^2-2ex\cos\alpha.}
\]
Cette méthode est rapide mais exige d'identifier sans ambiguïté le
triangle et l'angle opposé au côté considéré.

\subsection{Méthode 4 : orthogonalité}

Lorsqu'une géométrie impose deux directions perpendiculaires, on peut
écrire :
\[
\vec u\cdot\vec v=0.
\]
Cette méthode est particulièrement efficace pour certains renvois
d'angle, mécanismes à cardan et dispositifs de type Sinusmatic.

\begin{TLThreeApp}{Échelle EPAS}
Avec $\overline{O_0A}=a$, $\overline{AC}=c$,
$\overline{O_0B}=b$ et $\overline{BC}=r$, la géométrie conduit à :
\[
\overrightarrow{O_0A}+\overrightarrow{AC}
=\overrightarrow{O_0B}+\overrightarrow{BC}.
\]
En utilisant le triangle $O_0BC$ et les positions de $A$ et $C$ dans
le repère de l'échelle, établir $r=f(\theta)$, puis calculer
$|r(50^\circ)-r(10^\circ)|$ pour $a=2\ \mathrm m$,
$b=3\ \mathrm m$ et $c=2{,}6\ \mathrm m$.
\end{TLThreeApp}

\section{Loi entrée-sortie en vitesse}
\label{sec:les-vitesse}

\begin{TLThreeDef}{Définition}
La loi entrée-sortie en vitesse est la relation entre la dérivée du
paramètre de sortie et celle du paramètre d'entrée :
\[
\dot q_s=J(q_e,q_s)\,\dot q_e.
\]
Le coefficient $J$ est le rapport instantané de transmission. Il peut
dépendre de la configuration.
\end{TLThreeDef}

Deux démarches sont possibles :
\begin{itemize}
  \item dériver la loi en position déjà établie ;
  \item écrire directement une fermeture cinématique à l'aide des
  torseurs cinématiques ou des vitesses de points.
\end{itemize}

\begin{center}
\small
\begin{tabularx}{\linewidth}{>{\bfseries}p{3.5cm}X}
\toprule
Démarche & Enchaînement \\
\midrule
Fermeture géométrique & loi en position $\rightarrow$ dérivation temporelle $\rightarrow$ loi en vitesse \\
Fermeture cinématique & composition des vitesses ou des torseurs $\rightarrow$ élimination des paramètres internes \\
\bottomrule
\end{tabularx}
\end{center}

Rappels de dérivation :
\[
\frac{\mathrm d}{\mathrm dt}\sin\alpha
=\dot\alpha\cos\alpha,
\qquad
\frac{\mathrm d}{\mathrm dt}\cos\alpha
=-\dot\alpha\sin\alpha.
\]

Pour la loi implicite du bielle-manivelle :
\[
L^2=x^2+e^2-2ex\cos\alpha,
\]
la dérivation donne :
\[
0=2x\dot x-2e\dot x\cos\alpha
+2ex\dot\alpha\sin\alpha.
\]
Ainsi :
\[
\boxed{\frac{\dot x}{\dot\alpha}
=-\frac{ex\sin\alpha}{x-e\cos\alpha}.}
\]

\begin{TLThreeMethod}{Contrôler une loi entrée-sortie}
\begin{enumerate}
  \item Vérifier les unités : un rapport rotation--rotation est sans dimension ; un rapport translation--rotation est une longueur.
  \item Tester une position de référence facile à interpréter.
  \item Vérifier le signe avec les sens positifs du paramétrage.
  \item Repérer les configurations singulières où le dénominateur s'annule.
  \item Comparer l'évolution obtenue avec une animation ou une construction géométrique.
\end{enumerate}
\end{TLThreeMethod}

\section{Synthèse pour les mécanismes non linéaires}
\label{sec:les-synthese}

\begin{center}
\small
\begin{tabularx}{\linewidth}{>{\bfseries}p{3.2cm}X X}
\toprule
Structure & Modèle recherché & Outil privilégié \\
\midrule
Chaîne ouverte & position de l'effecteur & Chasles, projections, modèle direct \\
Chaîne ouverte, position imposée & paramètres articulaires & modèle inverse, résolution analytique ou numérique \\
Chaîne fermée plane & relation entre deux paramètres & fermeture géométrique, Al-Kashi, orthogonalité \\
Loi en vitesse & rapport instantané & dérivation ou fermeture cinématique \\
Transmetteur usuel & rapport constant ou loi simple & roulement sans glissement, engrènement, hélicoïdale \\
\bottomrule
\end{tabularx}
\end{center}

\section{Sources de la partie lois entrée-sortie}
\label{sec:les-sources}

Cette partie reprend le document de cours original, des ressources
pédagogiques de collègues de l'UPSTI et les systèmes issus de sujets de
concours mentionnés dans les exercices.
